\documentclass{article}

\usepackage{amsmath,amssymb}
\usepackage{xurl}
\usepackage[hidelinks]{hyperref}

% Shared commands for notes in docs/paper-gaps/.

\DeclareUnicodeCharacter{2080}{\ifmmode{}_{0}\else\textsubscript{0}\fi}
\DeclareUnicodeCharacter{2081}{\ifmmode{}_{1}\else\textsubscript{1}\fi}
\DeclareUnicodeCharacter{2082}{\ifmmode{}_{2}\else\textsubscript{2}\fi}
\DeclareUnicodeCharacter{2099}{\ifmmode{}_{n}\else\textsubscript{n}\fi}

\newcommand{\Lean}{\textsc{Lean}}
\ExplSyntaxOn
\NewDocumentCommand{\leanid}{m}{
  \ifmmode
    \textsf{\detokenize{#1}}
  \else
    \group_begin:
    \urlstyle{sf}
    \tl_set:Nn \l_tmpa_tl {#1}
    \tl_replace_all:Nnn \l_tmpa_tl {\_} {_}
    \exp_args:NV \path \l_tmpa_tl
    \group_end:
  \fi
}
\ExplSyntaxOff
\providecommand{\tr}{\operatorname{tr}}
\newcommand{\tnleanIssueBase}{https://github.com/LionSR/TNLean/issues/}
\newcommand{\tnleanPrBase}{https://github.com/LionSR/TNLean/pull/}
\newcommand{\tnleanDocBase}{https://sirui-lu.com/TNLean/docs/}
\newcommand{\ghissue}[1]{\href{\tnleanIssueBase#1}{GitHub issue \##1}}
\newcommand{\ghpr}[1]{\href{\tnleanPrBase#1}{PR \##1}}
\newcommand{\doclink}[1]{\href{\tnleanDocBase#1.html}{\path{#1}}}

% Dirac notation and the identity operator, as in blueprint/src/macros/common.tex.
% \providecommand keeps note-local definitions authoritative.
\usepackage{dsfont}
\providecommand{\ket}[1]{|#1\rangle}
\providecommand{\bra}[1]{\langle #1|}
\providecommand{\braket}[2]{\langle #1 | #2 \rangle}
\providecommand{\ketbra}[2]{|#1\rangle\!\langle #2|}
\providecommand{\Id}{\mathds{1}}
\providecommand{\one}{\mathds{1}}
\providecommand{\C}{\mathbb{C}}
\providecommand{\MN}[1]{M_{#1}(\C)}
\providecommand{\MD}{M_D(\C)}

% External referencing for paper sources.  The build helper writes sanitized
% auxiliary files under build/paper-aux/ and excludes duplicated source labels.
\usepackage{xr}

% Import a generated paper auxiliary file with a caller-supplied label prefix.
% The candidate paths support builds from docs/paper-gaps/, the repository root,
% and build/paper-aux/ itself.
\newcommand{\importpaperaux}[2]{%
  \IfFileExists{../../build/paper-aux/#2.aux}{%
    \externaldocument[#1]{../../build/paper-aux/#2}%
  }{%
    \IfFileExists{build/paper-aux/#2.aux}{%
      \externaldocument[#1]{build/paper-aux/#2}%
    }{%
      \IfFileExists{#2.aux}{%
        \externaldocument[#1]{#2}%
      }{}%
    }%
  }%
}

% General external reference macros
\newcommand{\paperref}[2]{\ref{#1-#2}}
\newcommand{\papereqref}[2]{(\ref{#1-#2})}

% CPSV16 source (1606.00608 / MPDO-22-12-17-2.tex) external document and shortcuts
\importpaperaux{cpsv16-}{1606.00608/MPDO-22-12-17-2-tnlean}
\newcommand{\cpsvref}[1]{\paperref{cpsv16}{#1}}
\newcommand{\cpsveqref}[1]{\papereqref{cpsv16}{#1}}

% Verdict marker: \gapnote{<kind>}{<status>}. Kind is the policy
% classification (clarification, local-correction, scope-restriction,
% unfaithful, false-source, open-gap); status is open, wip (elimination
% underway), resolved, or historical. Machine-read by the site index;
% prints a verdict line.
\newcommand{\gapnote}[2]{\par\noindent\textbf{Verdict:} #1 (#2).\par}


% Fallbacks keep this author document compilable when it is built outside its
% source directory and a different command.tex is found first.
\providecommand{\leanid}[1]{\textsf{\detokenize{#1}}}
\providecommand{\doclink}[1]{%
    \href{https://sirui-lu.com/TNLean/docs/#1.html}{\path{#1}}}
\providecommand{\gapnote}[2]{%
    \par\noindent\textbf{Verdict:} #1 (#2).\par}

\title{The Left Action Vectors of the CZX Product States}
\author{}
\date{2026-09-25}

\begin{document}
\maketitle
\gapnote{local-correction}{resolved}

The second computation of the CZX anomaly in
Ref.~\cite{GarreRubioSchuch2024DomainWall} obtains $\omega=-1$ from the
action tensors of the decorated CZX matrix product unitary $U$ on the
product states $|0\rangle^{\otimes N}$ and $|1\rangle^{\otimes N}$, which $U$
exchanges, through the ratio $L_0/L_1=-1$.\footnote{Source traceability:
\path{Papers/2405.00439/MPU-DW.tex}, lines 1246--1339; the printed vectors are
at lines 1272--1273 and 1300--1301.} The conclusion is correct, but the
printed left action vectors are not.

\section{The tensor and the action}

The tensor of $U$, with physical output $i$, input $j$, left bond $l$ and
right bond $r$, is
\begin{align}
    T^{ij}_{lr}
     & =\delta_{i,1-j}\,\hat H_{l,1-j}\,\delta_{r,1-j}\,(-1)^{1-j},
    \qquad
    \hat H=\begin{pmatrix}1&1\\1&-1\end{pmatrix}.
    \label{eq:gs24_czx_mpu_tensor}
\end{align}
The reading~\eqref{eq:gs24_czx_mpu_tensor} reproduces the printed fusion tensors
$\hat V=\langle1|\otimes\langle0|$ and $V=-|\hat-\rangle\otimes|\hat+\rangle$,
with $|\hat\pm\rangle=|0\rangle\pm|1\rangle$, and the sign $-1$ of their
associator.\footnote{Formal declarations:
\leanid{CZXCompression.czxDecoratedTensor},
\leanid{CZXCompression.czxDecorated_associator_eq_neg_one}.} The action
of $T$ on the product state $|s\rangle$ has the single nonzero letter
\begin{align}
    B_0^{1} & =T^{10}=-|\hat-\rangle\langle1|, &
    B_1^{0} & =T^{01}=|\hat+\rangle\langle0|.
    \label{eq:gs24_czx_action_letters}
\end{align}

\section{The printed vectors}

The source draws, for three sites, the decomposition
$B^{a}B^{b}B^{c}=B^{a}\,w\,\delta_{b,1-s}\,v\,B^{c}$ of the action on
$|s\rangle^{\otimes3}$ by the letters~\eqref{eq:gs24_czx_action_letters}, with
right vector $w$ after the first site and left vector $v$ before the third.
It prints $w=|\hat+\rangle$, $v=\langle\hat+|$
for $s=0$, and $w=-|\hat+\rangle$, $v=\langle\hat+|$ for $s=1$. For $s=0$ these
give $vw=2$ and $vB_0^{1}w=\langle\hat+|\hat-\rangle\cdot(\ldots)=0$, so the
right-hand side of the drawn decomposition vanishes although
$B_0^1B_0^1B_0^1\neq0$. For $s=1$ they give $vw=-2$. No choice among the four
readings of the diagram (the Hadamard on the left or right bond, the physical
input at the bottom or top), and no normalization factor, makes
$v=\langle\hat+|$ work.

\section{The correction}

Keeping the printed right vectors, the conditions $vw=1$,
$vB^{1-s}_sw=1$ and $B^{1-s}_swvB^{1-s}_s=(B^{1-s}_s)^2$ force
\begin{align}
    v & =\langle1| \quad (s=0), &
    v & \in\{(c,d):c+d=-1\},\ \text{for instance }-\langle0| \quad (s=1).
    \label{eq:gs24_czx_corrected_left_vectors}
\end{align}
With the choices~\eqref{eq:gs24_czx_corrected_left_vectors}, $v=\langle1|$ and
$v=-\langle0|$, both three-site decompositions hold and $(v,w)$
are reductions of the action onto the exchanged state. The resulting
L-symbols are $L^0_{s,s}=1$ and $L^1_{s,s}=-1$; the source's values
$L_i=(-1)^{i+1}$ differ from these by an action-tensor gauge, and the
gauge-invariant ratio $L_0/L_1=-1=\omega(s,s,s)$ of line 1335 is
unchanged.\footnote{Formal declarations:
\leanid{CZXCompression.czxBlockActionData},
\leanid{CZXCompression.czx_omega_eq_lSymbol_ratio},
\leanid{CZXCompression.czx_lSymbol_ratio_eq_neg_one} in
\doclink{TNLean/MPS/Symmetry/MPOSymmetry/CZXPermutedBlocks}.}

\textbf{Verdict: resolved local fix.} The left action vectors at lines 1272
and 1300 should read $\langle1|$ and $-\langle0|$; the anomaly computed from
them is the printed $\omega=-1$.

\bibliographystyle{alpha}
\bibliography{references}

\end{document}
